% Panel letters set by the manuscript layout, outside the artwork (RUNBOOK R5 and section 14:
% figures carry no text; panel letters sit outside the axes or the drawing).
% Layout geometry only, not prose: this file is not read by lint-numbers.
%
% \nacrePanelLetter{<offset>}{<letter>} sets a bold lower-case letter at <offset> from the left
% edge of the line, without taking width, so that several letters share one line. The letters of one figure sit in a
% box of the full line width, so a \centering figure does not move them to the middle of the line.
\newcommand{\nacrePanelLetter}[2]{\makebox[0pt][l]{\hspace*{#1}\textbf{\textsf{#2}}}}
%
% Fig 2 (figures/tikz/fig2.tex): the artwork is 17.1 cm wide and its four panels start at
% x = 0.05, 4.3, 8.55 and 12.8 cm. Each letter sits over the left edge of its panel, at the
% same fraction of the printed width; tests/preprint/test_preprint_model.py checks these
% numbers against the TikZ source.
\newcommand{\nacreFigTwoPanelLetters}{%
  \par\noindent
  \makebox[\linewidth][l]{%
  \nacrePanelLetter{\dimexpr\linewidth*5/1710\relax}{a}%
  \nacrePanelLetter{\dimexpr\linewidth*430/1710\relax}{b}%
  \nacrePanelLetter{\dimexpr\linewidth*855/1710\relax}{c}%
  \nacrePanelLetter{\dimexpr\linewidth*1280/1710\relax}{d}}%
  \par\nointerlineskip\vspace{1pt}\noindent}
%
% Fig 1 (figures/out/fig1.pdf, hero render): two separate single-phase panels in one 4039 x 1480 px
% image; panel (a) spans 0 to 977 and panel (b) starts at 1023 of 2000 horizontal units. Each letter
% sits above the left edge of its panel, outside the artwork.
\newcommand{\nacreFigOnePanelLetters}{%
  \par\noindent
  \makebox[\linewidth][l]{%
  \nacrePanelLetter{\dimexpr\linewidth*5/2000\relax}{a}%
  \nacrePanelLetter{\dimexpr\linewidth*1028/2000\relax}{b}}%
  \par\nointerlineskip\vspace{1pt}\noindent}
